% 使用 xelatex 编译
\documentclass[UTF8,a4paper,twoside,zihao=-4]{ctexrep}
%%%% 文档类设置 %%%%%%%%
\ctexset{
	contentsname = {目\quad 录},
	chapter = {
		beforeskip = 0pt,
		afterskip = 20pt,
		name = {第,章},
		nameformat = \zihao{3}\bfseries,
		number = \Chinese{chapter},
		numberformat = \zihao{3}\bfseries,
		titleformat = \zihao{3}\bfseries,
		fixskip = true,
		%afterindent = false
		        },
	section = {
		beforeskip = \ccwd,
		afterskip =1ex plus 0.2ex,
		format = \raggedright,
		nameformat = \zihao{4}\bfseries,
		numberformat = \zihao{4}\bfseries,
		titleformat = \zihao{4}\bfseries,
		%afterindent = false
			},
	subsection = {
		beforeskip =\ccwd,
		afterskip = 0.5ex,
		format = \raggedright,
		nameformat = \zihao{-4}\bfseries,
		numberformat = \zihao{-4}\bfseries,
		titleformat = \zihao{-4}\bfseries
			}
	    }

%%%% 宏包 %%%%%%%%%%%
\usepackage{amsmath,amssymb,amsfonts,mathrsfs}
\usepackage{graphicx}
\usepackage{tikz}
\usepackage{ulem}
\usepackage{pdfpages}
\usepackage{lipsum}
\usepackage{booktabs}
\usepackage{enumerate }
\usepackage[colorlinks,linkcolor=blue,citecolor=blue,bookmarks=true,bookmarksnumbered=true]{hyperref}
\usepackage[round]{natbib}
%%%%% 定理环境 %%%%%%%%%
\usepackage[thmmarks]{ntheorem}
{
  \theoremstyle{nonumberplain}
  \theoremheaderfont{\indent\bfseries}
  \theorembodyfont{\normalfont}
  \theoremsymbol{\ensuremath{\Box}}
  \newtheorem{proof}{证明}
}
{
  \theoremheaderfont{\indent\bfseries}
  \theorembodyfont{\normalfont}
  \newtheorem{theorem}{定理}[chapter]
  \newtheorem{lemma}{引理}[chapter]
  \newtheorem{definition}{定义}[chapter]
}
%%%%%% 图表标题设置 %%%%%%%%%%%%%%%
\usepackage{caption}
\DeclareCaptionLabelFormat{mylabel}{{\xeCJKsetup{CJKecglue={\hskip 0pt}}#1#2}}
\captionsetup{
		font= small,
		labelfont = bf,
		labelsep = space,
		labelformat = mylabel
			}
\makeatletter
\renewcommand{\thefigure}{\ifnum \c@chapter>\z@ \thechapter-\fi \@arabic\c@figure}
\renewcommand{\thetable}{\ifnum \c@chapter>\z@ \thechapter-\fi \@arabic\c@table}
\makeatother

%%%% 页面设置 %%%%%%%%%%%%%%%%%%%%%%
\usepackage[bindingoffset=.5cm,centering,includeheadfoot,margin=2.5cm,headsep=1em]{geometry}
\setlength\parskip{0pt}
\usepackage{fancyhdr}\pagestyle{fancy}
\renewcommand\chaptermark[1]{%
	\markboth{\CTEXthechapter\quad #1}{}}
\renewcommand\sectionmark[1]{%
	\markright{\CTEXthesection\quad #1}}
\fancyhf{}
\fancyhead[LE,RO]{\thepage}
\fancyhead[RE]{\textsl{\nouppercase{\leftmark}}}
\fancyhead[LO]{\textsl{\nouppercase{\rightmark}}}
%\renewcommand{\headrulewidth}{0pt}
	
	
%%%% 字体 %%%%%%%%%
% xelatex 默认字体为Latin Modern
%%%%% 自定义命令 %%%%%%%%%%%%
\newcommand{\nn}{\mathbf{N}^\ast}
\renewcommand{\leq}{\leqslant}
\renewcommand{\geq}{\geqslant}
\renewcommand{\ref}{\autoref}
\renewcommand{\epsilon}{\varepsilon}
%%%%%%%%%%%%%%%%%%%%%%
%%%%%%%%%%%%%%%%%%%%%%
\begin{document}
\bibliographystyle{plainnat}
%%%%%%%%% 封面 %%%%%%%%%%%%
%\includepdf{cover/cover.pdf}
%%%%%%% 中英文摘要 %%%%%%%%%%%%%%%%
\pagenumbering{roman}
% acknowledgement

\clearpage
\thispagestyle{plain}
\phantomsection
\addcontentsline{toc}{chapter}{Acknowledgment}

\centerline{\zihao{3}\bfseries Acknowledgment}

\linespread{1.4}\zihao{-4}
\bigskip

I would like to thank my mentor Dr.\,Xu Xiaodong for his support and patience throughout the whole writing and editing process of this thesis. My gratefulness also goes to Liu Yunhan and Zhong Yuling who helped me to solve crucial problems in experiment design and implementation.

The thesis is typeset with \LaTeX. When preparing the format template, I got much help from numerous dedicated members of different \LaTeX\ communities, to whom I express my dearest thankfulness.
%%% 中文摘要
\clearpage
\thispagestyle{plain}
\phantomsection
\addcontentsline{toc}{chapter}{摘\quad 要}

\centerline{\zihao{3}\heiti 摘\quad 要}

\linespread{1.4}\zihao{-4} \bigskip

本文探讨了英语和法语非母语者的语言处理过程中，焦点结构和代词回指之间的关系。首先我们从语法和语用角度，梳理了关于焦点作用和代词回指机制的研究现状。接着通过设计自定步速阅读实验，我们发现以分裂句为形式的焦点结构不一定能提升语言片段的显著性，并使其更可能成为代词回指的对象，这与前人的研究结果是相符的。在我们的测试中，法语主语位置的焦点和英语宾语位置的焦点反应时间都较短，但两种语言的主语焦点都导致了更长的回指反应时间。另外，我们还发现焦点和回指对象的一致性并不能提升焦点位置的显著性。因此本文认为，英语和法语中是否存在主语或宾语回指偏向的问题，比当前已有研究结论更加复杂。
\bigskip

\noindent{\zihao{4}\heiti 关键词:}
焦点作用, 代词回指, 自定步速阅读, 英语, 法语

% Abstract
\clearpage
\thispagestyle{plain}
\phantomsection
\addcontentsline{toc}{chapter}{Abstract}

\centerline{\zihao{3}\bfseries Abstract}

\linespread{1.4}\zihao{-4}
\bigskip

This thesis explores the relationship between focus structure and pronoun resolution among non-native speakers of English and French. Firstly we reviewed the existing literature on the mechanism of focus effect and pronoun resolution. Then through a self-paced reading test, we find that focus, in the form of cleft structure does not necessarily increase the salience of a informational unit, thus may not in some cases make it a preferred antecedent for pronoun resolution. This result is line with previous researches on this topic. In our experiment, We also find that focused subject in French and focused object in English are processed faster, but focused subjects in both languages leads to longer response time of anaphora. Furthermore, our research also shows that the congruence between anaphora and focus does not make the latter more accessible. In this regard, we argue that the problem of whether there is subject or object preference in English and French is more complicated than the results of current studies.

\bigskip
\noindent\textbf{\zihao{4} Keywords:} 
focus effect, pronoun resolution, self-paced reading, English, French


\tableofcontents\newpage\mbox{}\thispagestyle{empty}\newpage
\clearpage\pagenumbering{arabic}
\chapter{}
\chapter{}
\chapter{}
\begin{thebibliography}{99}
\addcontentsline{toc}{chapter}{参考文献}

\bibitem[常庚哲, 史济怀(2016)]{CS2016}
常庚哲, 史济怀, 2016, 《数学分析教程(上册)第三版》[M]. 合肥: 中国科学技术大学出版社. 
\end{thebibliography}

\end{document}